\documentclass[12pt,a4paper]{article}

% ==================================================
% 1. CONFIGURATION
% ==================================================

% Enable / disable title page
\newif\ifshowtitlepage
\showtitlepagetrue
% \showtitlepagefalse

% Document information
\newcommand{\worktitle}{Running a Choreography}
\newcommand{\worksubtitle}{NAO Robot Exercise System}
\newcommand{\teamname}{NAOMI}
\newcommand{\subject}{NAO Robot Exercise System}

\newcommand{\university}{Slovak University of Technology in Bratislava}
\newcommand{\faculty}{Faculty of Informatics and Information Technologies}

\newcommand{\academicyear}{2026/2027}
\newcommand{\location}{Bratislava}
\newcommand{\documentyear}{2026/2027}

% Team leaders
\newcommand{\leaderone}{Ing. Jakub Perdek, PhD.}
\newcommand{\leadertwo}{ doc. Ing. Fedor Lehocki PhD.}

% Team members
\newcommand{\memberone}{Bc. Akhmadinurov Amal}
\newcommand{\membertwo}{Bc. Danilov Denis}
\newcommand{\memberthree}{Bc. Dmitriev Anton}
\newcommand{\memberfour}{Bc. Gainetdinov Ruslan}
\newcommand{\memberfive}{Bc. Kuzin Timofei}
\newcommand{\membersix}{Bc. Tikhoblazhenko Alena}
\newcommand{\memberseven}{Bc. Zimnitski Aliaksei}


% ==================================================
% 2. PACKAGES
% ==================================================

\usepackage[
    a4paper,
    top=3cm,
    bottom=2.7cm,
    left=2.8cm,
    right=2.8cm,
    headheight=45pt,
    headsep=20pt
]{geometry}

\usepackage[T1]{fontenc}
\usepackage[utf8]{inputenc}
\usepackage[english]{babel}
\usepackage{lmodern}

\usepackage{graphicx}
\graphicspath{{../../assets/}{assets/}}
\usepackage[table]{xcolor}
\usepackage{tikz}
\usetikzlibrary{arrows.meta,positioning}
\usepackage{listings}
\usepackage{enumitem}
\setlist[itemize]{itemsep=1pt,parsep=0pt,topsep=2pt}
\lstset{basicstyle=\ttfamily\scriptsize,breaklines=true,frame=single,columns=fullflexible,keepspaces=true}
\usepackage{fancyhdr}
\usepackage{titlesec}
\usepackage{setspace}
\usepackage{lastpage}
\usepackage{array}
\usepackage{float}
\usepackage{caption}
\usepackage{booktabs}
\usepackage{tabularx}
\usepackage{lipsum}
\usepackage{csquotes}

% Hyperlinks
\usepackage[hidelinks]{hyperref}


% ==================================================
% 3. COLORS
% ==================================================

\definecolor{primary}{HTML}{000000}
\definecolor{secondary}{HTML}{000000}
\definecolor{muted}{HTML}{000000}
\definecolor{linecolor}{HTML}{000000}
\definecolor{naomiTableHeader}{HTML}{17445A}
\definecolor{naomiTableStripe}{HTML}{EDF7FA}
\definecolor{naomiTableRule}{HTML}{66BBD0}

% Ragged-right columns keep narrow table cells readable.
\newcolumntype{L}[1]{>{\raggedright\arraybackslash}p{#1}}
\newcolumntype{Y}{>{\raggedright\arraybackslash}X}


% ==================================================
% 4. TYPOGRAPHY
% ==================================================

\setstretch{1.2}

\setlength{\parindent}{0pt}
\setlength{\parskip}{7pt}

% Section
\titleformat{\section}
    {\Large\bfseries\color{primary}}
    {\thesection}
    {0.8em}
    {}

% Subsection
\titleformat{\subsection}
    {\large\bfseries\color{primary}}
    {\thesubsection}
    {0.8em}
    {}

% Subsubsection
\titleformat{\subsubsection}
    {\normalsize\bfseries\color{primary}}
    {\thesubsubsection}
    {0.8em}
    {}

\titlespacing*{\section}{0pt}{20pt}{10pt}
\titlespacing*{\subsection}{0pt}{14pt}{7pt}
\titlespacing*{\subsubsection}{0pt}{12pt}{5pt}

% Captions
\captionsetup{
    font=small,
    labelfont=bf,
    justification=centering
}

% ==================================================
% 5. HEADER AND FOOTER
% ==================================================

\pagestyle{fancy}
\fancyhf{}

% Header logo
\fancyhead[L]{
    \includegraphics[
        width=4cm,
        height=1.15cm,
        keepaspectratio
    ]{horizontal.png}
}

% Header information
\fancyhead[R]{
    \begin{tabular}{r}
        {\footnotesize\bfseries \worksubtitle}\\[3pt]
        {\scriptsize Academic Year \academicyear}
    \end{tabular}
}

% Header line
\renewcommand{\headrulewidth}{0.7pt}

% Footer left
\fancyfoot[L]{
    {\scriptsize
    \teamname\enspace|\enspace\subject}
}

% Footer right
\fancyfoot[R]{
    {\scriptsize
    \thepage\enspace/\enspace\pageref{LastPage}}
}

% Footer line
\renewcommand{\footrulewidth}{0.4pt}

% Plain pages
\fancypagestyle{plain}{
    \fancyhf{}

    \fancyhead[L]{
        \includegraphics[
            width=4cm,
            height=1.15cm,
            keepaspectratio
        ]{horizontal.png}
    }

    \fancyhead[R]{
        \begin{tabular}{r}
            {\footnotesize\bfseries \worksubtitle}\\[3pt]
            {\scriptsize Academic Year \academicyear}
        \end{tabular}
    }

    \fancyfoot[L]{
        {\scriptsize
        \teamname\enspace|\enspace\subject}
    }

    \fancyfoot[R]{
        {\scriptsize
        \thepage\enspace/\enspace\pageref{LastPage}}
    }

    \renewcommand{\headrulewidth}{0.7pt}
    \renewcommand{\footrulewidth}{0.4pt}
}




% ==================================================
% 6. TITLE PAGE
% ==================================================

\newcommand{\makemytitlepage}{
    \begin{titlepage}
        \thispagestyle{empty}
        \centering

        % ==========================================
        % UNIVERSITY LOGO
        % ==========================================

        \vspace*{0.3cm}

        \includegraphics[
            width=6cm,
            height=5cm,
            keepaspectratio
        ]{vertical.png}

        \vspace{0.5cm}

        {\large\bfseries
        \university\par}

        \vspace{0.2cm}

        {\normalsize
        \faculty\par}

        % ==========================================
        % PROJECT TITLE
        % ==========================================

        \vspace{1.7cm}

        {\small\MakeUppercase{\worksubtitle}\par}

        \vspace{0.5cm}

        \begin{minipage}{0.95\textwidth}
            \centering

            {\fontsize{25}{32}\selectfont\bfseries
            \worktitle\par}
        \end{minipage}

        % ==========================================
        % TEAM INFORMATION
        % ==========================================

        \vfill

        \begin{minipage}{0.9\textwidth}
            \centering
            \begin{minipage}[t]{0.46\linewidth}
                \centering
                {\normalsize\bfseries Team Leaders\par}
                \vspace{0.25cm}
                {\small
                \leaderone\par
                \leadertwo\par}
            \end{minipage}%
            \hfill
            \begin{minipage}[t]{0.48\linewidth}
                \centering
                {\normalsize\bfseries Team Members\par}
                \vspace{0.25cm}
                {\small
                \memberone\par
                \membertwo\par
                \memberthree\par
                \memberfour\par
                \memberfive\par
                \membersix\par
                \memberseven\par}
            \end{minipage}
        \end{minipage}

        % ==========================================
        % BOTTOM INFORMATION
        % ==========================================

        \vfill

        {\small
        \location\enspace\textbullet\enspace\academicyear\par}

    \end{titlepage}
}

% ==================================================
% 7. IMAGE PLACEHOLDER
% ==================================================

% Usage:
% \imageplaceholder{Text}{Caption}

\newcommand{\imageplaceholder}[2]{

    \begin{figure}[H]
        \centering

        \begingroup
        \setlength{\fboxsep}{0pt}

        \fbox{
            \begin{minipage}[c][5cm][c]{0.85\textwidth}

                \centering

                {\Large\bfseries IMAGE}

                \vspace{0.3cm}

                {\normalsize #1}

            \end{minipage}
        }

        \endgroup

        \caption{#2}

    \end{figure}
}


% ==================================================
% 8. REAL IMAGE COMMAND
% ==================================================

% Usage:
% \documentimage{path}{caption}{label}

\newcommand{\documentimage}[3]{

    \begin{figure}[htbp]
        \centering

        \includegraphics[
            width=0.85\textwidth,
            keepaspectratio
        ]{#1}

        \caption{#2}
        \label{#3}

    \end{figure}
}


% ==================================================
% 9. TABLES
% ==================================================

% Usage:
% \begin{naomitable}{Caption}{label}
%     \begin{tabularx}{\textwidth}{L{2.8cm} Y L{2.6cm}}
%         \rowcolor{naomiTableHeader}
%         \textcolor{white}{\bfseries Name} & ... \\
%         ...
%         \arrayrulecolor{naomiTableRule}\bottomrule[1pt]
%     \end{tabularx}
% \end{naomitable}

\newenvironment{naomitable}[2]{%
    \begin{table}[H]
        \centering
        \caption{#1}
        \label{#2}
        \begingroup
        \small
        \renewcommand{\arraystretch}{1.4}
        \setlength{\tabcolsep}{7pt}
        \rowcolors{2}{naomiTableStripe}{white}
}{%
        \endgroup
    \end{table}
}

% Usage:
% \tableplaceholder{Caption}{label}

\newcommand{\tableplaceholder}[2]{
    \begin{naomitable}{#1}{#2}
        \begin{tabularx}{\textwidth}{L{2.8cm} Y L{2.6cm}}
            \rowcolor{naomiTableHeader}
            \textcolor{white}{\bfseries Name} &
            \textcolor{white}{\bfseries Description} &
            \textcolor{white}{\bfseries Status} \\
            Item A & Sample description & Completed \\
            Item B & Sample description & In Progress \\
            Item C & Sample description & Planned \\
            \arrayrulecolor{naomiTableRule}\bottomrule[1pt]
        \end{tabularx}
    \end{naomitable}
}


% ==================================================
% 10. DOCUMENT
% ==================================================

\begin{document}
\ifshowtitlepage
    \makemytitlepage
\fi

\section{Purpose}

This document describes the runtime API for a multi-state choreography. The
choreography graph is stored by \texttt{exerciseConfigService}; each state
references a saved exercise. \texttt{translator} loads that graph at start and
controls playback on one robot. The graph editor and automatic assessment of a
participant's movement are outside this runtime contract.

A run has its own \texttt{run\_id}. A successful performance ends in
\texttt{waiting}; only then can the client send a transition signal or move to
another state without playback. The run remains queryable after it stops.

\section{Run Lifecycle}

\begin{naomitable}{Run statuses}{tab:statuses}
\begin{tabularx}{\textwidth}{L{2.5cm} Y}
\rowcolor{naomiTableHeader}
\textcolor{white}{\bfseries Status} & \textcolor{white}{\bfseries Meaning} \\
\texttt{performing} & The initial or target exercise is being played. \\
\texttt{waiting} & Playback is complete; signal or goto can be accepted. \\
\texttt{stopped} & Playback has ended or was stopped; signals are rejected. \\
\arrayrulecolor{naomiTableRule}\bottomrule[1pt]
\end{tabularx}
\end{naomitable}

The start and signal endpoints return \texttt{202} after scheduling work in a
background thread. The exercise may still be playing when the response arrives.
The client polls the run endpoint for its status and result. Voice lines precede
poses in each stored frame; the saved frame sequence and repetition count are
used. The graph is a snapshot loaded at start, while exercise data is fetched
when a state is performed.

\section{HTTP Contract}

\begin{naomitable}{Runtime endpoints on translator, port 7000}{tab:run-api}
\begin{tabularx}{\textwidth}{L{1.6cm} L{5.1cm} Y}
\rowcolor{naomiTableHeader}
\textcolor{white}{\bfseries Method} & \textcolor{white}{\bfseries Path} & \textcolor{white}{\bfseries Effect} \\
POST & \texttt{/choreography/<id>/run} & Start and schedule the initial state; no body. \\
GET & \texttt{/run/current} & Get the active run; \texttt{404} if none. \\
GET & \texttt{/run/<run\_id>} & Get state, status, transitions and playback outcome. \\
POST & \texttt{/run/<run\_id>/signal} & Send \texttt{\{"transition\_id":"go"\}} and schedule the target exercise. \\
POST & \texttt{/run/<run\_id>/goto} & Send \texttt{\{"state\_id":"<exercise\_id>"\}}; change state without robot commands. \\
POST & \texttt{/run/<run\_id>/stop} & Stop between robot commands; reject later signals. \\
\arrayrulecolor{naomiTableRule}\bottomrule[1pt]
\end{tabularx}
\end{naomitable}

The \texttt{<id>} in the start path is the stored choreography ID. The
\texttt{<run\_id>} returned by start identifies this particular execution.
The example transition ID \texttt{go} is an ordinary label chosen by the
choreography author, not a reserved command. An exercise ID in a request must
identify a state already present in the loaded graph.

\subsection{Run response}

All successful runtime endpoints return the same run representation. For
example, after the initial exercise completes:

\begin{lstlisting}
{
  "run_id": "<run_id>",
  "choreography_id": "<choreography_id>",
  "state_id": "<first_exercise_id>",
  "status": "waiting",
  "outgoing_transitions": [
    {"id": "go", "from": "<first_exercise_id>",
     "to": "<second_exercise_id>", "label": "Continue"}
  ],
  "playback": {"exercise_id": "<first_exercise_id>",
               "repetitions": 1, "frames_per_rep": 1,
               "total_poses_sent": 2, "results": ["..."]},
  "error": null
}
\end{lstlisting}

During \texttt{performing}, \texttt{playback} is the previous result or
\texttt{null}. If playback fails, \texttt{error} explains the failure and
\texttt{playback} contains the commands attempted so far, when available.
\texttt{total\_poses\_sent} is inherited from the standalone exercise API;
it counts result entries, including voice commands. The listed outgoing
transitions are those whose \texttt{from} equals the current \texttt{state\_id}.

\subsection{Errors}

\begin{naomitable}{HTTP errors}{tab:errors}
\begin{tabularx}{\textwidth}{L{1.6cm} Y}
\rowcolor{naomiTableHeader}
\textcolor{white}{\bfseries Status} & \textcolor{white}{\bfseries Meaning} \\
\texttt{400} & Invalid choreography ID or missing JSON field. \\
\texttt{404} & Unknown choreography, run, state or transition; no active run for \texttt{/run/current}. \\
\texttt{409} & Another run is active; the run is not waiting; the transition starts elsewhere; or stop is repeated. \\
\texttt{502} & The choreography service is unavailable or returns an invalid graph. \\
\arrayrulecolor{naomiTableRule}\bottomrule[1pt]
\end{tabularx}
\end{naomitable}

An empty or unavailable exercise is discovered by the background worker:
start or signal has already returned \texttt{202}, and the run endpoint then
reports the failure in \texttt{error}. Failed initial playback changes the
status to \texttt{stopped}. Failed signal playback returns to \texttt{waiting}
with the source state still current. Commands already accepted by the robot
cannot be undone, so retrying a signal can repeat them.

\section{Stop and Concurrency}

Only one choreography run is active per translator process. Another start
returns \texttt{409}, even for another choreography ID. A stop request marks
the run \texttt{stopped} immediately and prevents further commands between
frames or repetitions. A robot call already in progress may finish; the robot
remains reserved until the worker exits. The standalone exercise endpoint
\texttt{POST /exercise/<id>/run} retains its existing response contract and
shares the playback lock.

Run data and locks are held in translator memory. A restart loses them;
multiple translator processes do not coordinate their runs. Other robot
control routes do not share the playback lock. The production NAO adapter
and physical robot need separate testing before a live demonstration.

\section{Routing and Scope}

Vite, the container Nginx configuration and the Kubernetes web ConfigMap
route \texttt{/run/} and \texttt{/choreography/<id>/run} to translator.
Choreography CRUD remains in \texttt{exerciseConfigService}. The runtime
selects only an explicitly supplied transition ID. Automatic selection,
movement assessment, conditions and repetition counters belong to later work.

\section{Client Walkthrough}

Use an existing choreography ID when starting. Keep the returned
\texttt{run\_id} for later calls. The following requests show the sequence;
the placeholders must be replaced with actual IDs.

\begin{lstlisting}
POST /choreography/<choreography_id>/run
GET  /run/<run_id>                  # poll until waiting
POST /run/<run_id>/signal
Content-Type: application/json

{"transition_id":"go"}
GET  /run/<run_id>                  # poll until waiting
POST /run/<run_id>/goto
Content-Type: application/json

{"state_id":"<exercise_id>"}
POST /run/<run_id>/stop
\end{lstlisting}

The client can use \texttt{GET /run/current} to discover the active run
without remembering its ID. It must use the returned
\texttt{outgoing\_transitions} to present valid choices. A transition whose
\texttt{from} differs from the current state is rejected without a robot call.

\section{Verification Boundary}

Offline tests use Flask's test client, a fake exercise service and fake robot
calls. They cover playback order, transitions, direct moves, stopping during a
long command, concurrent starts and robot failure. A Choregraphe virtual NAO
can support a manual walkthrough; it does not replace verification through the
production adapter or a separately coordinated physical robot test.

\end{document}
